\documentclass[11pt]{article}
\usepackage[margin=1in]{geometry}
\usepackage{enumitem}
\usepackage{tabularx}
\usepackage{xltabular}
\usepackage{amsmath}
\usepackage{amssymb}
\usepackage{multirow}
\usepackage{array}

\setlength{\parindent}{0pt}
\renewcommand{\arraystretch}{1.3}

\begin{document}

\textbf{DOK 3 TASK QUESTIONS} \hfill \textit{Reverse-engineer a polynomial from its graph}

\vspace{0.5em}
\noindent\textit{DOK arc for the day: gradual-release Launch (I-do $\rightarrow$ you-do $\rightarrow$ I-check $\rightarrow$ we-do-together) gives students a worked example to lean on, THEN the DOK-3 Mystery Graph hits in Explore, then DOK-2 consolidation, then a DOK-1 cool-down at Exit. The DOK-3 task is loaded near the front of Explore (not buried at the end) so the reasoning happens while energy is high.}

\vspace{1em}

\noindent \textbullet \ \textbf{Entry / Access (build understanding, low floor)}
\begin{itemize}[label=$\bullet$, leftmargin=0.5in]
    \item Where does this graph cross the $x$-axis? Where does it touch without crossing?
    \item What is the lowest possible degree this graph could be? Why?
    \item What does the end behavior tell you about the leading coefficient?
    \item If a graph touches but does not cross at $x=2$, what kind of factor must $(x-2)$ be?
\end{itemize}

\vspace{1em}

\noindent \textbullet \ \textbf{Developing (connect representations, reasoning)}
\begin{itemize}[label=$\bullet$, leftmargin=0.5in]
    \item How do the zeros and their multiplicities together determine the shape of the graph near each zero?
    \item Why does an even multiplicity produce a touch and an odd multiplicity produce a cross?
    \item Two students propose different factored forms for the same graph. How would you decide which (if either) is correct?
\end{itemize}

\vspace{1em}

\noindent \textbullet \ \textbf{Core Task (DOK 3 -- reasoning \& justification, FRONT-LOADED)}
\begin{itemize}[label=$\bullet$, leftmargin=0.5in]
    \item Given only a graph (zeros at $x=-1$ touch, $x=3$ cross; ends down-down), construct a possible factored form. Justify each factor.
    \item Could there be more than one polynomial with this graph? What is the same about all of them, and what can vary?
    \item Defend your leading coefficient sign using end behavior, not by plugging in numbers.
    \item A classmate writes $f(x) = (x+1)(x-3)$ for the same graph. Build a counterexample showing why their answer fails the touch-at-$-1$ requirement.
\end{itemize}

\vspace{1em}

\noindent \textbullet \ \textbf{Extension (higher-level thinking, generalization)}
\begin{itemize}[label=$\bullet$, leftmargin=0.5in]
    \item Generalize: from any graph with labeled touches, crosses, and end behavior, can you always recover a factored form (up to a constant)?
    \item How would the answer change if the touch at $x=-1$ were a multiplicity-4 touch instead of multiplicity-2? What would the graph look like differently?
\end{itemize}

\vspace{1em}

\noindent \textbullet \ \textbf{Discussion / Argumentation prompts (to support collaboration)}
\begin{itemize}[label=$\bullet$, leftmargin=0.5in]
    \item Do you agree with your partner's factored form? Where exactly do you disagree?
    \item What evidence from the graph best supports your factor for $x=-1$?
    \item Can there be more than one reasonable approach? Compare them side-by-side.
\end{itemize}

\newpage

% --- Lesson Plan Start ---

\noindent
\begin{tabularx}{\textwidth}{|X|}
\hline
\textbf{Rationale:} Students learn to read multiplicity off a graph through gradual release (I-do, you-do, I-check, we-do-together), then apply the rule by reverse-engineering a polynomial from a graph in a DOK-3 task at the start of Explore. The worked-example scaffold de-risks the reasoning task while preserving the productive-struggle moment. \\
\textbf{Materials:} L35\_P2 student packet, projected ``Mystery Graph'' slide, sentence-frame card, Desmos available for the share-out. \\
\hline
\end{tabularx}

\vspace{-1.2pt}
\noindent
\begin{tabularx}{\textwidth}{|>{\hsize=0.8\hsize}X|>{\hsize=0.8\hsize}X|>{\hsize=1.4\hsize}X|}
\hline
\textbf{Date:} 5/8/26 & \textbf{Grade Level:} 11 & \textbf{Subject:} Algebra 2 -- Lesson 3-5, Period 2 (Multiplicity + Factored-Form Graphing) \\
\hline
\end{tabularx}

\vspace{-1.2pt}
\noindent
\begin{tabularx}{\textwidth}{|X|}
\hline
\textbf{Common Core State Standards:}
\begin{itemize}[label=$\bullet$, leftmargin=*]
    \item \textbf{HSA.APR.B.3:} Identify zeros of polynomials when suitable factorizations are available, and use the zeros to construct a rough graph of the function defined by the polynomial.
    \item \textbf{HSF.IF.C.7.C:} Graph polynomial functions, identifying zeros when suitable factorizations are available, and showing end behavior.
\end{itemize}

\textbf{Knowledge Needed:}
\begin{itemize}[label=$\bullet$, leftmargin=*, nosep]
    \item Factored form of a polynomial $\rightarrow$ zeros connection
    \item End behavior depends on degree parity and leading-coefficient sign
    \item Tier 2 vocabulary: zero, root, multiplicity, factor, end behavior, degree, leading coefficient
\end{itemize}
\vspace{0.5em}

\textbf{Skills Needed:}
\begin{itemize}[label=$\bullet$, leftmargin=*, nosep]
    \item Read $x$-intercepts from a graph
    \item Distinguish a touch from a cross visually
    \item Convert factored form to expanded form (and back) for verification
\end{itemize}
\vspace{0.5em}

\textbf{Standards for Mathematical Practice}
\begin{itemize}[label=, leftmargin=*]
    \item 1. Make sense of problems and persevere in solving them
    \item 3. Construct viable arguments and critique the reasoning of others
    \item 7. Look for and make use of structure
\end{itemize} \\
\hline
\textbf{Unit Goals:} Multiplicity is the bridge between symbolic factored form and graph behavior. Students who own this connection can later read rational functions, exponential models, and trig graphs structurally rather than by memorization. \\
\hline
\end{tabularx}

\vspace{-1.2pt}
\noindent
\begin{tabularx}{\textwidth}{|>{\hsize=1.1\hsize}X|>{\hsize=0.9\hsize}X|}
\hline
\textbf{Content Objective:} \newline\newline
1. Construct a factored polynomial from a graph by reasoning about zeros, multiplicity, and end behavior. \newline\newline
2. Justify why even multiplicity produces a touch and odd multiplicity produces a cross at a zero. &
\textbf{Criteria for Success/Formative Assessment:}
\begin{itemize}[label=$\bullet$, leftmargin=*, nosep]
    \item \textbf{TPS check-ins during Launch and Explore (4--6 circulation laps)}
    \item \textbf{Exit ticket: vocabulary + 1 quick reverse-engineer}
\end{itemize}
\vspace{0.5em}
\textbf{Self-Reflection Assessment:}
\begin{itemize}[label=$\bullet$, leftmargin=*, nosep]
    \item \textbf{``The biggest thing I learned today was \_\_\_\_.''}
\end{itemize}
\vspace{0.5em}
\textbf{Teacher Assessment}
\begin{itemize}[label=$\bullet$, leftmargin=*, nosep]
    \item \textbf{Exit-ticket review before P3 (Storage Box) on Mon 5/11}
\end{itemize} \\
\hline
\end{tabularx}

\newpage

\noindent
\begin{tabularx}{\textwidth}{|X|X|}
\hline
\multicolumn{2}{|l|}{\textbf{Language Objective:}} \\
\multicolumn{2}{|l|}{
\begin{minipage}{\linewidth}
\begin{itemize}[label=$\bullet$, leftmargin=*]
    \item Students will justify a chosen factored form from a graph, verbally and in writing, using sentence frames and structured think-pair-share.
\end{itemize}
\end{minipage}
} \\
\hline
\textbf{Modifications or Support for IEP students.}
\begin{itemize}[label=$\bullet$, leftmargin=*, nosep]
    \item Rephrasing/CFU of expectations
    \item Pre-printed graph copy taped into packet for annotation
    \item Extended TPS think-time on cue
\end{itemize} &
\textbf{Supports for ELL students}
\begin{itemize}[label=$\bullet$, leftmargin=*, nosep]
    \item Vocabulary card: zero, root, multiplicity, touch, cross, end behavior
    \item Sentence starters during TPS:
\end{itemize}
``The factor for $x=\_\_\_\_$ must be \_\_\_\_ because the graph \_\_\_\_'' \newline
``The leading coefficient is \_\_\_\_ because the ends go \_\_\_\_'' \newline
``I notice \_\_\_\_, so I think \_\_\_\_''
\begin{itemize}[label=$\bullet$, leftmargin=*, nosep]
    \item Two-reads strategy on the Mystery Graph prompt
\end{itemize} \\
\hline
\end{tabularx}

\vspace{1.5em}
\begin{center}
\textbf{Lesson Sequence}
\end{center}

\noindent
\begin{xltabular}{\textwidth}{|>{\raggedright\arraybackslash\hsize=0.7\hsize}X|
                               >{\raggedright\arraybackslash\hsize=1.2\hsize}X|
                               >{\raggedright\arraybackslash\hsize=1.2\hsize}X|
                               >{\raggedright\arraybackslash\hsize=0.9\hsize}X|}
\hline
\textbf{Lesson Part} & \textbf{Teachers will...} & \textbf{Students will...} & \textit{Questions to consider...} \\
\hline
\endhead

% ROW 1: Do Now
\textbf{Do Now (5 min)} \newline\newline Quick warm-up, NOT the DOK-3 hook &
Project 3 factored polynomials. For each, list the zeros (no graphing yet).
\scriptsize
\begin{enumerate}[label=\arabic*., leftmargin=*, nosep]
    \item $f(x) = (x-2)(x+5)$ \quad [zeros: $2, -5$]
    \item $g(x) = (x-1)^2(x+3)$ \quad [zeros: $1, -3$]
    \item $h(x) = x(x-4)(x+2)$ \quad [zeros: $0, 4, -2$]
\end{enumerate} \normalsize
\textbf{Adult role:} Spot-check; flag \#2 (the squared factor) for the Launch hook. &
Read each polynomial; list zeros; circle anything unusual. &
How can you engage students in the lesson? \newline\newline
How are all adults intentionally planned for? \\
\hline

% ROW 2: Launch -- gradual release
\textbf{Launch (15 min):} \newline \textbf{\textit{Gradual release}}
\begin{itemize}[label=$\bullet$, leftmargin=*, nosep]
    \item I do
    \item You do
    \item I check
    \item We do together
\end{itemize} &
\textbf{I do (4 min).} Project $f(x) = (x-2)^2(x+3)$. Think aloud while sketching: ``Zero at $x=2$, factor squared $\rightarrow$ even multiplicity $\rightarrow$ touches. Zero at $x=-3$, factor to the first power $\rightarrow$ odd multiplicity $\rightarrow$ crosses. Degree 3, leading coefficient $+$, so end behavior down-up.'' Name the touch/cross rule explicitly on the board. \newline\newline
\textbf{You do (4 min).} Project $g(x) = (x+1)(x-4)^2$. Students sketch independently in their packet using the rule. \newline\newline
\textbf{I check (3 min).} Circulate. Pull two student sketches to project (one correct, one with a common error like swapping touch/cross). Name the misconception out loud, fix it as a class. \newline\newline
\textbf{We do together (4 min).} Project $h(x) = -x(x+2)^2(x-1)$ (one harder: leading coefficient negative, three zeros). Whole-class call-and-response: ``Where are the zeros?'' ``Which one touches?'' ``Why does it open downward?'' Teacher writes; students copy. \newline\newline
Read content + language objectives aloud at the end. \newline\newline
\textbf{Adult role:} Model precisely in I-do; circulate aggressively in You-do; name errors without judgment in I-check; conduct We-do as a chorus. &
\textbf{I do:} Watch and annotate the example in the packet; copy the touch/cross rule. \newline\newline
\textbf{You do:} Independently sketch $g(x) = (x+1)(x-4)^2$ using the rule. \newline\newline
\textbf{I check:} Compare with the projected exemplar; revise; participate in the misconception fix. \newline\newline
\textbf{We do together:} Call out answers chorally for the harder example; copy the final sketch. &
What will you explicitly model and how? \newline\newline
How will you know if students are ready to work alone? \newline\newline
How are all adults intentionally planned for? \\
\hline

% ROW 3: Explore -- DOK-3 task + DOK-2 consolidation
\textbf{Explore (25 min):} \newline \textbf{\textit{DOK-3 task, then DOK-2 consolidation}}
\begin{itemize}[label=$\bullet$, leftmargin=*, nosep]
    \item Mystery Graph (DOK-3, 10 min)
    \item Mixed graph $\leftrightarrow$ factored practice (DOK-2, 15 min)
    \item Students doing the heavy lifting
\end{itemize} &
\textbf{Mystery Graph (10 min, TPS).} Project: zeros at $x=-1$ (touch), $x=3$ (cross); ends both pointing \emph{down}. \newline\newline
Pose: ``Build a possible factored form for this graph. Justify every piece -- the factors, their powers, and the leading coefficient.'' Note for students: this is harder than Launch -- the ends-down case forces you to reason about leading coefficient AND degree together. \newline\newline
\textbf{Start TPS -- 3 min independent, 3 min partner share, 4 min whole-group.} \newline\newline
Circulate; tag 2--3 different responses. In whole-group share, press: ``Why must the leading coefficient be negative?'' ``Could the degree be 4 instead of 3? What would change?'' \newline\newline
\textbf{DOK-2 consolidation (15 min, TPS).} Hand out 3 practice items from L35 P2 bank: \newline
-- Item 1: factored $\rightarrow$ graph \newline
-- Item 2: factored $\rightarrow$ graph (with a multiplicity-3 factor) \newline
-- Item 3: graph $\rightarrow$ factored (lower-stakes than Mystery Graph) \newline\newline
\textbf{Questions while circulating:}
\begin{itemize}[label=$\bullet$, leftmargin=*, nosep]
    \item Where does the graph touch vs.\ cross? Why?
    \item What does end behavior tell you about degree and sign?
    \item How would you check your factored form is right?
\end{itemize}
\textbf{Adult role:} 4--6 circulation laps; pull a struggling pair to a re-teach huddle if needed; tag exemplars for Share. &
\textbf{Mystery Graph:} TPS -- propose a factored form; defend each factor and the leading coefficient using the sentence frame. \newline\newline
\textbf{DOK-2 items:} work in pairs; record reasoning, not just answers; verify with Desmos at the end of each item. &
What will students do for most of the session? What will they practice? How? \newline\newline
What type of academic conversations? \newline\newline
How will you be interacting with them? \\
\hline

% ROW 4: Share/Summary
\textbf{Share/Summary (5 min)}
\begin{itemize}[label=$\bullet$, leftmargin=*, nosep]
    \item Synthesis
    \item Return to objectives
    \item Name the rule again
\end{itemize} &
Re-state the multiplicity rule using student language from Launch. Press: ``Where in today's items did the rule save you time?'' \newline\newline
Project objectives. Ask: ``Which one did you hit? Where did you struggle?'' &
Share one item where the rule helped; agree/disagree/add on; reflect on objectives. &
How will students express what they have learned? \\
\hline

% ROW 5: Exit -- DOK-1 COOL-DOWN
\textbf{Exit Ticket (5 min)} \newline \textbf{\textit{DOK-1 cool-down}} &
Project exit prompt (low-stakes by design):
\begin{enumerate}[label=\arabic*., leftmargin=*, nosep]
    \item A factor of $(x-4)^2$ produces a \underline{\hspace{2em}} at $x=4$ (touch / cross).
    \item A factor of $(x+1)$ produces a \underline{\hspace{2em}} at $x=-1$ (touch / cross).
    \item One sentence: ``The biggest thing I learned today was \_\_\_\_.''
\end{enumerate} &
Complete exit ticket independently; turn in on the way out. &
How will students show what they've learned based on the objective? \\
\hline

\end{xltabular}

\end{document}
